\ticket{23}{Современные нейросетевые языковые модели: особенности архитектуры и обучения, области применения. Понятие переноса обучения (Transfer Learning) и тонкой настройки (Fine-Tuning) на примере модели BERT}

\begin{examplan}
    \item Описать особенности современных нейросетевых языковых моделей.
    \item Раскрыть архитектуру Transformer.
    \item Объяснить предобучение, transfer learning и fine-tuning.
    \item Показать тонкую настройку на примере BERT.
\end{examplan}

\subsection*{Короткая версия для письменного ответа}

Современные нейросетевые языковые модели используют большие корпуса и архитектуру Transformer, чтобы строить контекстуальные представления текста. В отличие от n-грамм, они способны учитывать дальние зависимости и переносить знания на разные NLP-задачи.

\subsection*{Transformer}

Ключевые компоненты:
\begin{itemize}
    \item \textbf{self-attention}: каждое слово учитывает другие слова последовательности;
    \item \textbf{multi-head attention}: несколько голов внимания выделяют разные типы зависимостей;
    \item \textbf{позиционные кодировки}: добавляют информацию о порядке слов;
    \item \textbf{feed-forward слои}: нелинейная обработка представлений;
    \item \textbf{residual connections и layer normalization}: стабилизируют обучение глубоких моделей.
\end{itemize}

Типы архитектур:
\begin{itemize}
    \item encoder-only: BERT, RoBERTa; задачи понимания текста;
    \item decoder-only: GPT; генерация текста;
    \item encoder-decoder: T5, BART, машинный перевод и суммаризация.
\end{itemize}

\subsection*{Обучение}

Обычно используется два этапа.

\textbf{Предобучение} на больших неразмеченных корпусах с задачами самообучения:
\begin{itemize}
    \item Masked Language Modeling: предсказать замаскированные токены;
    \item Causal Language Modeling: предсказать следующий токен по предыдущим;
    \item дополнительные задачи, например NSP в исходном BERT.
\end{itemize}

\textbf{Дообучение} адаптирует предобученную модель к конкретной задаче на размеченных данных.

\subsection*{Transfer Learning}

\textbf{Transfer learning} --- перенос знаний, полученных на одной задаче, на другую связанную задачу. В NLP модель сначала учит общие закономерности языка на больших корпусах, а затем эти веса используются как начальная точка для классификации, NER, QA, анализа тональности и других задач.

Преимущества:
\begin{itemize}
    \item меньше потребность в размеченных данных;
    \item выше качество на целевых задачах;
    \item быстрее обучение по сравнению с обучением с нуля.
\end{itemize}

\subsection*{BERT и fine-tuning}

\textbf{BERT} --- encoder-only модель Transformer, строящая двунаправленные контекстные представления. Она предобучалась на Masked Language Modeling и Next Sentence Prediction.

Тонкая настройка BERT:
\begin{enumerate}
    \item выбрать предобученную модель;
    \item подготовить вход: \texttt{[CLS]}, токены, \texttt{[SEP]}, attention mask;
    \item добавить задачно-специфичную голову;
    \item дообучить модель на размеченных данных с малой скоростью обучения;
    \item оценить качество на тестовой выборке.
\end{enumerate}

Примеры:
\begin{itemize}
    \item классификация текста: использовать представление \texttt{[CLS]} и softmax-слой;
    \item NER: классифицировать каждый токен;
    \item QA: предсказать начало и конец ответа в тексте.
\end{itemize}

Варианты: полное дообучение всех параметров или использование BERT как экстрактора признаков с замороженными слоями.

\begin{important}
    \item Transformer основан на механизме внимания, а не на рекуррентном проходе.
    \item Современные модели дают контекстуальные эмбеддинги.
    \item Предобучение использует большие неразмеченные корпуса.
    \item Transfer learning переносит языковые знания на прикладные задачи.
    \item Fine-tuning BERT обычно добавляет небольшую голову и дообучает модель под задачу.
\end{important}
